% TEMPLATE-MILC-v3.tex - plantilla maestra UNICA de las guias MILC v3.
%
% La usan las tres familias de guias, cada una con su driver:
%   · Grados 6-11  -> scripts/build-guias-g11.py
%   · Bebras       -> scripts/build-guias-bebras.py
%   · Semillero    -> scripts/build-guias-semillero.py
%
% Antes habia tres copias casi identicas (~400 lineas iguales, 6 distintas):
% cada mejora habia que aplicarla tres veces, y en la practica no se hacia
% (el motor de recursos vivia solo en dos de ellas). Lo que varia entre
% familias esta parametrizado en 6 placeholders que cada driver llena
% (se nombran aqui SIN su sintaxis real: el builder reemplaza por texto plano,
% tambien dentro de los comentarios, y un valor multilinea rompe el preambulo):
%
%   HEADER_IZQ        encabezado izquierdo de pagina
%   HEADER_DER        encabezado derecho de pagina
%   PORTADA_ETIQUETA  rotulo sobre el numero grande (GRADO/MOMENTO/MODULO)
%   PORTADA_SERIE     linea de serie de la portada
%   PORTADA_ID        identificador de la guia en la portada
%   PORTADA_CIRCULO   el circulito de grado (°), o vacio si no aplica
%
% El resto de claves las documenta content/guias/_SCHEMA.md. El script valida
% que no queden placeholders sin reemplazar antes de escribir el archivo final.

\documentclass[11pt,letterpaper]{article}
\usepackage[margin=1.75cm]{geometry}
\usepackage[table]{xcolor}
\usepackage{fontspec}
\usepackage[spanish]{babel}
\usepackage{tikz}
% Librerías TikZ para los diagramas de las guías (bloque `recursos:`):
%  · babel  → babel en español vuelve activos caracteres como «>», y TikZ los usa
%             en sus opciones (>=stealth, ->). Sin esta librería, cualquier
%             diagrama con flechas revienta con «\language@active@arg> has an
%             extra }». Debe ir ANTES de que se dibuje nada.
%  · positioning        → sintaxis «below left=2pt and 3pt».
%  · decorations.pathreplacing → llaves ({brace}) para señalar rangos.
\usetikzlibrary{babel,positioning,decorations.pathreplacing}
\usepackage{graphicx}
% Circuitos: el trabajo de electrónica se hace hoy con micro:bit y sus sensores
% integrados (MakeCode, sin cableado), así que no se necesitan esquemáticos.
% Si algún día entran montajes con protoboard, basta descomentar esta línea:
% los diagramas .tex/.tikz declarados en `recursos:` ya se incrustan solos.
% \usepackage{circuitikz}
\usepackage[most]{tcolorbox}
\usepackage{tabularx}
\usepackage{array}
\usepackage{fancyhdr}
\usepackage{titlesec}
\usepackage{ragged2e}
\usepackage{enumitem}
\usepackage{microtype}

% Helvetica Neue no trae la flecha U+2192, asi que cada «→» escrito literalmente
% en el YAML se compilaba a NADA, en silencio (462 flechas en 61 guias: rutas del
% tipo «paleta → tipografia → lockup» salian rotas). Mapeamos el caracter a su
% version matematica, que si tiene glifo. Asi el autor puede seguir escribiendo
% «→» comodamente en el YAML.
\usepackage{newunicodechar}
\newunicodechar{→}{\ensuremath{\rightarrow}}

\setmainfont{Helvetica Neue}
\setsansfont{Helvetica Neue}
\setlength{\parindent}{0pt}
\setlength{\parskip}{6pt}
\hyphenpenalty=200
\exhyphenpenalty=200
\pretolerance=400
\tolerance=2000
\setlength{\emergencystretch}{1em}
\renewcommand{\arraystretch}{1.25}
\setlist{leftmargin=5mm,itemsep=1mm,topsep=1mm}
\newcolumntype{Y}{>{\RaggedRight\arraybackslash}X}

\definecolor{milcMagenta}{HTML}{E6007E}
\definecolor{milcVino}{HTML}{5A0038}
\definecolor{milcTurquesa}{HTML}{0093A5}
\definecolor{milcVerde}{HTML}{58A923}
\definecolor{milcMostaza}{HTML}{E5B400}
\definecolor{milcOcre}{HTML}{B86600}
\definecolor{milcNegro}{HTML}{241816}
\definecolor{milcGris}{HTML}{F7F7F4}
\definecolor{milcLinea}{HTML}{E8E2DD}
\definecolor{milcGradePrimary}{HTML}{7C3AED}
\definecolor{milcGradeDark}{HTML}{4C1D95}
\definecolor{milcGradeSoft}{HTML}{EDE9FE}
\definecolor{milcGradeAccent}{HTML}{CFE7FF}
\definecolor{milcGradeText}{HTML}{FFFFFF}
\color{milcNegro}

\pagestyle{fancy}
\fancyhf{}
\lhead{\small Tecnología e Informática · Grado 9}
\rhead{\small Guía 20 · MILC v3}
\cfoot{\small\thepage}
\renewcommand{\headrulewidth}{0pt}

\newtcolorbox{titlebox}[1]{
  enhanced,colback=#1,colframe=#1,arc=7mm,boxrule=0pt,
  left=7mm,right=7mm,top=3mm,bottom=3mm,
  before skip=8pt,after skip=8pt,
  fontupper=\sffamily\bfseries\Large\color{white}
}

\newtcolorbox{softbox}[3]{
  enhanced,colback=#2,colframe=#1,borderline west={4pt}{0pt}{#1},
  arc=2mm,boxrule=.4pt,left=4mm,right=4mm,top=3mm,bottom=3mm,
  before skip=6pt,after skip=8pt,title={#3},coltitle=white,
  colbacktitle=#1,fonttitle=\sffamily\bfseries,fontupper=\RaggedRight,
  attach boxed title to top left={xshift=3mm,yshift=-2mm},
  boxed title style={arc=2mm, boxrule=0pt}
}

\newtcolorbox{drawbox}[2]{
  enhanced,colback=white,colframe=#1,arc=4mm,boxrule=1pt,
  left=5mm,right=5mm,top=4mm,bottom=4mm,before skip=8pt,after skip=8pt,
  title={#2},coltitle=white,colbacktitle=#1,fonttitle=\sffamily\bfseries,
  fontupper=\RaggedRight,
  attach boxed title to top left={xshift=4mm,yshift=-2mm},
  boxed title style={arc=3mm, boxrule=0pt}
}

\newtcolorbox{infoband}[1]{
  enhanced,colback=#1!8,colframe=#1!50,arc=2mm,boxrule=.5pt,
  left=4mm,right=4mm,top=2mm,bottom=2mm,before skip=4pt,after skip=4pt,
  fontupper=\small\RaggedRight
}

% Puente narrativo entre secciones (contrato editorial Conéctate).
% Da continuidad: "Acabas de X. Ahora vas a Y porque Z."
% Visualmente: banda izquierda en color del grado, texto en itálica pequeña.
\newtcolorbox{puentebox}{
  enhanced,colback=milcGradeSoft,colframe=milcGradeDark,
  borderline west={3pt}{0pt}{milcGradeDark},
  arc=0mm,boxrule=0pt,left=5mm,right=4mm,top=2.5mm,bottom=2.5mm,
  before skip=4pt,after skip=8pt,
  fontupper=\itshape\small\color{milcNegro}\RaggedRight
}
% La flecha va en modo matematico a proposito: Helvetica Neue NO tiene el glifo
% U+2192, asi que \textrightarrow se compilaba a NADA («Missing character: There
% is no → in font Helvetica Neue») y el puente salia sin flecha. En modo
% matematico la flecha viene de la fuente matematica y si se dibuja.
\newcommand{\puente}[1]{%
  \begin{puentebox}%
    \textcolor{milcGradeDark}{$\rightarrow$}\hspace{2mm}#1%
  \end{puentebox}%
}

\newcommand{\linead}[1]{\noindent\color{milcLinea}\rule{#1}{0.5pt}\color{milcNegro}}
\newcommand{\respuesta}[1]{\par\vspace{2mm}\foreach \n in {1,...,#1}{\noindent\color{milcLinea}\rule{\linewidth}{0.5pt}\par\vspace{5mm}}\color{milcNegro}}
\newcommand{\checkbox}{\textcolor{milcGradeDark}{\fbox{\phantom{x}}}\hspace{5pt}}

\newcommand{\phasecard}[4]{
\begin{tcolorbox}[enhanced,colback=#3,colframe=#2,arc=3mm,boxrule=.5pt,height=3.05cm,valign=center,halign=center,left=1.5mm,right=1.5mm,top=2mm,bottom=2mm]
{\sffamily\bfseries\fontsize{22}{24}\selectfont\textcolor{#2}{#1}}\par
{\sffamily\bfseries\small #4}
\end{tcolorbox}
}

% ── Recursos visuales de la guía ──────────────────────────────────────
% Declarados en el bloque `recursos:` del YAML y emitidos por el builder.
% \guiaFigura   → imagen rasterizada o PDF (foto, carta celeste, montaje).
% \guiaDiagrama → fuente TikZ/circuitikz (.tex) incrustada como vector:
%                 es la vía para circuitos y esquemáticos editables.
% #1 = ruta relativa al .tex · #2 = pie (puede ir vacío)
\newcommand{\guiaFigura}[2]{%
\begin{center}
\includegraphics[width=0.88\linewidth,height=0.38\textheight,keepaspectratio]{#1}\\[1.5mm]
{\footnotesize\itshape\textcolor{milcNegro!75}{#2}}
\end{center}
}

\newcommand{\guiaDiagrama}[2]{%
\begin{center}
\input{#1}\\[1.5mm]
{\footnotesize\itshape\textcolor{milcNegro!75}{#2}}
\end{center}
}

\begin{document}
\thispagestyle{empty}
\begin{tikzpicture}[remember picture,overlay]
  \fill[white] (current page.south west) rectangle (current page.north east);
  \path[fill=milcGradePrimary,rounded corners=16mm]
    ([xshift=.55cm,yshift=-.55cm]current page.north west) rectangle
    ([xshift=-.55cm,yshift=3.65cm]current page.south east);

  \node[anchor=north west,text=milcGradeText] at ([xshift=2.0cm,yshift=-1.85cm]current page.north west)
    {\sffamily\bfseries\fontsize{14}{16}\selectfont GRADO};
  \node[anchor=north west,text=milcGradeText,opacity=.82] at ([xshift=2.0cm,yshift=-2.75cm]current page.north west)
    {\sffamily\bfseries\fontsize{9}{11}\selectfont SERIE GUÍAS · TECNOLOGÍA E INFORMÁTICA};

  \begin{scope}[shift={([xshift=-3.05cm,yshift=-2.55cm]current page.north east)}]
    \fill[white,opacity=.80] (-.62,-.52) rectangle (.62,.52);
    \draw[milcGradeText,line width=1pt] (-.62,-.52) rectangle (.62,.52);
    \fill[milcGradeDark] (-.42,-.38) rectangle (-.22,.06);
    \fill[milcGradeAccent] (-.10,-.38) rectangle (.10,.24);
    \fill[milcGradePrimary!60!black] (.22,-.38) rectangle (.42,.38);
  \end{scope}

  \node[anchor=west,text=milcGradeText] at ([xshift=1.9cm,yshift=-12.85cm]current page.north west)
    {\sffamily\bfseries\fontsize{124}{124}\selectfont 9};
  % El circulito de grado (°) solo aplica a las guias de grado («11°»); en
  % Bebras y Semillero el numero grande es un momento/modulo, no un grado.
  \draw[milcGradeText,line width=1.7pt]
    ([xshift=4.52cm,yshift=-11.68cm]current page.north west) circle (.13cm);

  \node[anchor=west,text=milcGradeText,text width=8.7cm,align=left] at ([xshift=10.35cm,yshift=-11.6cm]current page.north west)
    {
     {\sffamily\bfseries\fontsize{19}{22}\selectfont Cosecha P2 ---\\sustentación pública\\de la revista}\\[4mm]
     {\sffamily\bfseries\fontsize{23}{26}\selectfont Guía 20-9-TIC}\\[2mm]
     {\sffamily\fontsize{13}{16}\selectfont Período 2 · Diseño editorial digital}\\[5mm]
     {\sffamily\bfseries\fontsize{11}{13}\selectfont Institución Educativa Sor María Juliana}\\[1mm]
     {\sffamily\fontsize{10}{12}\selectfont Autor: Dr. Álvaro Cárdenas Orozco}
    };

  \path[fill=milcGradeSoft,rounded corners=7mm]
    ([xshift=1.35cm,yshift=.85cm]current page.south west) rectangle
    ([xshift=-1.35cm,yshift=4.25cm]current page.south east);
  \draw[milcGradeDark,line width=.65pt,rounded corners=7mm]
    ([xshift=1.35cm,yshift=.85cm]current page.south west) rectangle
    ([xshift=-1.35cm,yshift=4.25cm]current page.south east);
  \node[anchor=center,text=milcNegro,text width=17.0cm,align=center] at ([yshift=2.55cm]current page.south)
    {\sffamily\fontsize{10}{12}\selectfont Metodología MILC · Apertura ancestral · Pensamiento computacional · Triángulo Dussel-Estoicismo-Floridi\\
    \textcolor{milcGradeDark}{\bfseries Producto:} Exposición oral pública de la revista al curso (1 minuto por página, total 8 min), donde el estudiante defiende sus decisiones de diseño y recibe feedback estructurado de 3 compañeros};
\end{tikzpicture}
\mbox{}
\newpage

\section*{Apertura: Cosecha P2 --- sustentación pública de la revista}

\begin{softbox}{milcGradeDark}{milcGradeSoft}{Saber ancestral}
Los \textbf{palabreros Wayuu} de La Guajira no improvisan: cuando representan a una familia en un conflicto, ensayan el discurso durante días, lo someten al juicio de los mayores y solo entonces lo presentan en público. Equivocarse en una palabra puede romper un pacto entre clanes. En las \textbf{asambleas Páez (Nasa)} del Cauca, el joven que va a hablar por primera vez en público es presentado por un \emph{taita} mayor que avala su preparación. Esa misma tradición de \emph{sustentación con maestros y pares} sostiene la \textbf{tradición de los talleres editoriales colombianos}: en las imprentas de Cali, Cartago y el Eje Cafetero, el aprendiz tipógrafo nunca publicaba su primer pliego sin presentarlo al maestro y a los pares. Soportaba críticas amigas: ``este interlineado es muy apretado'', ``este titular pesa más de lo necesario'', ``aquí hay un viudo''. \textbf{Sustentar es someter el trabajo a una mirada externa antes de que el público lo lea.} Esa práctica antigua --- defender en voz alta lo que hiciste, escuchar lo que otros ven que tú no, ajustar --- es lo que separa el trabajo de aula del trabajo profesional.
\end{softbox}

\begin{softbox}{milcTurquesa}{milcGris}{Saber contemporáneo}
Hoy sustentar tu revista digital ante el curso es lo más cercano que tendrás --- a tu edad --- al \emph{review} profesional. En 8 minutos defiendes 8 páginas: decisión por decisión, sin esconder lo débil, sin inflar lo fuerte. Después recibes feedback estructurado de 3 compañeros. \textbf{No es prueba de nervios: es entrenamiento de oficio.} Cada presentación que hagas en tu vida adulta --- proyecto en universidad, propuesta a un cliente, pitch a un inversor --- depende de habilidades que se entrenan aquí.
\end{softbox}

\begin{softbox}{milcMostaza}{milcGris}{Pregunta-puente}
\textit{Si tuvieras 8 minutos para defender ante 30 personas el trabajo más importante de tu vida hasta hoy, ¿qué dirías? ¿Cómo demostrarías que no es solo ``una tarea''?}
\end{softbox}

\begin{softbox}{milcVerde}{milcGris}{Saber y saber hacer}
Al terminar podrás: (1) \textbf{analizar} tu propia revista identificando 3 fortalezas y 3 debilidades; (2) \textbf{evaluar} qué decisión de diseño defenderías con más convicción y cuál cambiarías; (3) \textbf{explicar} oralmente 8 páginas en 8 minutos con orden, criterio y voz propia; (4) \textbf{crear} estructura de sustentación + entrega oral pública.
\end{softbox}



\small
\begin{tabularx}{\textwidth}{>{\bfseries}p{3.0cm}Y}
\rowcolor{milcGradePrimary!12}Autor & Dr. Álvaro Cárdenas Orozco\\
DBA & Producción de piezas editoriales digitales con criterio estético y ético (MEN, T\&I)\\
Referentes & Tipografía colombiana · Diseño centrado en accesibilidad · Floridi\\
Producto final & Exposición oral pública de la revista al curso (1 minuto por página, total 8 min), donde el estudiante defiende sus decisiones de diseño y recibe feedback estructurado de 3 compañeros\\
\end{tabularx}
\normalsize

\puente{Antes de empezar, mira el viaje completo. Nueve sesiones te trajeron hasta aquí. Hoy haces la cosecha pública: defiendes lo que produjiste y aceptas mirada externa. Ese ritual cierra el periodo.}

\begin{titlebox}{milcGradeDark}
Ruta MILC: así vamos a aprender
\end{titlebox}
\begin{tabularx}{\textwidth}{>{\centering\arraybackslash}X>{\centering\arraybackslash}X>{\centering\arraybackslash}X>{\centering\arraybackslash}X}
\phasecard{1}{milcVerde}{milcGris}{Escucha\\[-1mm]\small Observo, escucho y pregunto.} &
\phasecard{2}{milcTurquesa}{milcGris}{Sistematización\\[-1mm]\small Pensamiento computacional.} &
\phasecard{3}{milcMagenta}{milcGris}{Praxis\\[-1mm]\small Construyo y aplico.} &
\phasecard{4}{milcVino}{milcGris}{Evaluación\\[-1mm]\small Reflexión liberadora.}
\end{tabularx}


\puente{Antes de sustentar tu propia revista, vamos a escuchar 3 sustentaciones reales (sea de TED, de pitchs cortos, de exposiciones académicas). Verás patrones --- las buenas presentaciones comparten estructura.}

\begin{titlebox}{milcVerde}
Fase 1 · Escucha
\end{titlebox}

\begin{softbox}{milcVerde}{milcGris}{Escena para escuchar}
\textbf{Actividad 1 · IDENTIFICA --- 3 sustentaciones reales} (15 min · individual). Busca \textbf{3 presentaciones} reales de máximo 10 minutos: (a) una \emph{charla TED corta} (puedes filtrar por duración), (b) un \emph{pitch de proyecto} de algún concurso o startup (busca ``pitch 5 min'' o ``pitch competition''), (c) una \emph{sustentación académica} (defensa de tesis o proyecto de aula). Para cada una: ¿cómo empieza?, ¿cómo desarrolla?, ¿cómo cierra?
\end{softbox}

\checkbox Vi 3 presentaciones reales de máximo 10 minutos cada una.\\
\checkbox Identifiqué cómo empieza, desarrolla y cierra cada presentación.\\
\checkbox Reconocí qué hace fuerte a las buenas (estructura clara, voz propia, ejemplos concretos).

\begin{infoband}{milcVerde}


\textbf{Lo que va al cuaderno (Actividad 1 · IDENTIFICA).} Título: «Actividad 1 --- 3 sustentaciones reales». \textbf{Formato:} tabla de 4 columnas (Presentación~|~Inicio~|~Desarrollo~|~Cierre), 3 filas. Al final 1 párrafo de 3 renglones sobre el patrón común de las buenas presentaciones. \textbf{Sabes que terminaste cuando} reconoces patrones que vas a aplicar a tu propia sustentación.
\end{infoband}

\puente{Ya viste presentaciones reales. Ahora vamos a entender los \textbf{4 elementos} de una sustentación efectiva: contexto (qué hiciste), proceso (cómo decidiste), defensa (por qué eso y no otra cosa), reconocimiento (qué cambiarías). Sin alguno, la sustentación queda débil.}

\begin{titlebox}{milcTurquesa}
Fase 2 · Sistematización con pensamiento computacional
\end{titlebox}

\begin{softbox}{milcTurquesa}{milcGris}{Cuatro pilares aplicados al tema}
Toda sustentación efectiva tiene 4 elementos. Vamos a descomponerlos con los pilares del pensamiento computacional para que tu defensa de 8 minutos sea sólida.
\end{softbox}

\small
\begin{tabularx}{\textwidth}{>{\bfseries\RaggedRight\arraybackslash}p{3.6cm}Y}
\rowcolor{milcTurquesa!90}\color{white}Pilar & \color{white}\bfseries Aplicación al tema\\
1. Descomposición & La sustentación se descompone en 4 elementos: (1) \textbf{contexto} (qué hiciste --- 1 min: tema de la revista + 8 páginas + URL), (2) \textbf{proceso} (cómo decidiste --- 2-3 min: grilla elegida, tipografía, paleta, ¿por qué?), (3) \textbf{defensa} (página por página --- 3-4 min: decisión más fuerte de cada una), (4) \textbf{reconocimiento} (qué cambiarías --- 1 min: honestidad sobre debilidades).\\
2. Reconocimiento de patrones & Toda sustentación fuerte sigue el patrón \textbf{``honestidad + orden''}: si sostienes una decisión con seguridad pero también reconoces lo que cambiarías, el jurado/audiencia confía más. Si solo defiendes sin matiz, suena arrogante. Si solo reconoces sin defender, suena inseguro. Las dos cosas juntas son madurez profesional.\\
3. Abstracción & Lo esencial cabe en una pregunta: \emph{¿puedo defender cada decisión visual con un argumento que un compañero entendería en 30 segundos?}. Si la respuesta es no, hay que ensayar.\\
4. Algoritmo conceptual & Pasos para preparar la sustentación: (1) ensaya en voz alta sola/o 2 veces midiendo el tiempo; (2) anota los puntos clave por página en una tarjeta (no leer la hoja); (3) prepara la URL en un dispositivo listo; (4) presenta sin esconder lo débil ni inflar lo fuerte.\\
\end{tabularx}
\normalsize

\begin{drawbox}{milcTurquesa}{Anatomía de una sustentación de 8 minutos}
\textbf{Minuto 1 · Contexto:} tema + público objetivo + URL. \\[1mm] \textbf{Minutos 2-3 · Proceso:} grilla, tipografía, paleta elegidos y por qué. \\[1mm] \textbf{Minutos 4-7 · Defensa:} pasa página por página, 30 segundos cada una, defendiendo decisión más fuerte. \\[1mm] \textbf{Minuto 8 · Reconocimiento:} qué cambiarías en una versión 2.0 y qué te llevas para el siguiente periodo.
\end{drawbox}

\begin{softbox}{milcMostaza}{milcGris}{Errores comunes y su corrección}
(1) Leer la hoja en lugar de exponer: pierde conexión con la audiencia. (2) Justificar todo sin reconocer debilidades: suena defensivo. (3) Solo enumerar lo que hizo sin explicar por qué: pierde la oportunidad de demostrar criterio. (4) Hablar más de 8 minutos: la disciplina del tiempo importa. (5) No ensayar: se nota inmediatamente.
\end{softbox}

\begin{infoband}{milcTurquesa}


\textbf{Lo que va al cuaderno (Actividad 2 · EXPLICA).} Título: «Actividad 2 --- Estructura de mi sustentación». \textbf{Formato:} 4 fichas, una por elemento (contexto, proceso, defensa, reconocimiento), cada una con: minutos asignados + 3-5 puntos clave que vas a decir. \textbf{Sabes que terminaste cuando} podrías hablar 8 min en voz alta sin notas adicionales --- esto es ensayo, no improvisación.
\end{infoband}

\puente{Tienes el método. Toca producir: estructura de sustentación + entrega oral pública de 8 minutos ante el curso. No es nota --- es ritual de cierre del periodo.}

\begin{titlebox}{milcMagenta}
Fase 3 · Praxis con pensamiento computacional
\end{titlebox}

\begin{softbox}{milcMagenta}{milcGris}{Construyendo el producto con los cuatro pilares}
\textbf{Actividad 3 · CREA --- Sustentación oral pública de tu revista} (45 min · individual + grupal). Hora de defender. Vas a presentar tu revista de 8 páginas al curso durante 8 minutos. Antes, vas a tener 10 minutos de preparación final. Después, vas a recibir feedback estructurado de 3 compañeros.
\end{softbox}

\small
\begin{tabularx}{\textwidth}{>{\bfseries\RaggedRight\arraybackslash}p{3.6cm}Y}
\rowcolor{milcMagenta!90}\color{white}Pilar & \color{white}\bfseries Aplicación al producto\\
Descomposición & Tu presentación se descompone en 3 fases: (1) preparación final (10 min antes de exponer): repaso de la tarjeta de puntos clave, prueba de la URL, ensayo mental rápido; (2) entrega oral (8 min frente al curso): sigue la estructura 1-2-3-4 (contexto/proceso/defensa/reconocimiento); (3) feedback (5 min): escuchas 3 observaciones de compañeros sin defenderte.\\
Patrones & Sigue el patrón \textbf{``segura sin arrogancia''}: defiende cada decisión, pero reconoce honestamente lo que cambiarías. La gente confía en quien dice ``esto lo haría mejor en v2'' tanto como en quien dice ``estoy orgulloso de esto y aquí está por qué''.\\
Abstracción & Cada minuto expresa una sola idea. No condenses contexto + proceso + defensa en 3 minutos --- pierde estructura. Respeta los tiempos.\\
Algoritmo de armado & Algoritmo de la presentación: (1) preséntate y di el tema en 30 segundos; (2) muestra URL pública; (3) explica decisiones de sistema visual en 2-3 minutos; (4) recorre las 8 páginas en 4 minutos defendiendo decisión más fuerte de cada una; (5) cierra con 1 minuto de qué cambiarías y qué te llevas. Al terminar, escucha 3 observaciones de compañeros sin interrumpir.\\
\end{tabularx}
\normalsize

\begin{drawbox}{milcMagenta}{Checklist antes de declarar la sustentación hecha}
\checkbox Estructura escrita en cuaderno (4 elementos con minutos asignados). \\[1mm] \checkbox Ensayé al menos 2 veces en voz alta midiendo el tiempo. \\[1mm] \checkbox URL pública de la revista lista y verificada. \\[1mm] \checkbox Defendí oralmente la revista durante 7-8 minutos. \\[1mm] \checkbox Reconocí al menos 1 debilidad honestamente. \\[1mm] \checkbox Escuché 3 observaciones de compañeros sin defenderme.
\end{drawbox}

\begin{softbox}{milcMagenta}{milcGris}{Plantilla de trabajo}
\textbf{Estructura de mi sustentación.} \textit{Las 4 partes con minutos asignados.} \\[1mm] \textbf{Tarjeta de puntos clave.} \textit{Lista breve por minuto, no para leer literal.} \\[1mm] \textbf{Feedback recibido.} \textit{Las 3 observaciones literales de compañeros.} \\[1mm] \textbf{Reflexión final.} \textit{Qué cambia mi visión del trabajo después de la sustentación.}
\end{softbox}

\begin{infoband}{milcMagenta}


\textbf{Lo que va al cuaderno (Actividad 3 · CREA).} Título: «Actividad 3 --- Mi sustentación pública». \textbf{Formato:} (1) estructura de 4 elementos con minutos, (2) tarjeta de puntos clave usada, (3) feedback literal de 3 compañeros, (4) reflexión final de 3 renglones. \textbf{Extensión:} 2 páginas de cuaderno. \textbf{Sabes que terminaste cuando} defendiste oralmente y escuchaste feedback sin defenderte.
\end{infoband}

\puente{Lo que sigue es el resumen de qué entregas hoy: estructura escrita + sustentación oral defendida. El criterio principal: que defiendas tu revista en voz alta sin titubear y reconozcas con honestidad lo que cambiarías.}

\begin{titlebox}{milcMostaza}
Producto de la sesión
\end{titlebox}
\textbf{Sustentación oral pública de la revista (8 min) + feedback estructurado recibido}.
\begin{softbox}{milcMostaza}{milcGris}{Criterios de calidad}
(1) Estructura escrita previa (4 elementos con minutos). (2) Ensayo previo de al menos 2 veces en voz alta. (3) Sustentación oral de 7-8 minutos (no menos, no más). (4) Defensa de decisiones con argumentos + reconocimiento honesto de debilidad. (5) Feedback literal de 3 compañeros transcrito. (6) Reflexión final sobre lo aprendido del proceso.
\end{softbox}

\puente{Antes de cerrar, mira el acto de sustentar desde las cinco dimensiones humanas. Defender en voz alta lo que hiciste es ciudadanía digital, soberanía personal, herencia de oficio. Vale la pena nombrar las tres.}

\begin{titlebox}{milcVino}
Fase 4 · Evaluación liberadora — cinco dimensiones
\end{titlebox}

\begin{softbox}{milcVino}{milcGris}{Sentido de la evaluación}
Evaluar no es solo revisar una nota. Es reconocer qué aprendí, qué cambió en mí, qué emociones manejé, cómo me relaciono con la ciudad y la comunidad, y qué le dejaré a quien viene detrás.
\end{softbox}

\textbf{Mi reflexión liberadora en cinco dimensiones.} Completa cada frase iniciada:

\small
\begin{tabularx}{\textwidth}{>{\bfseries\RaggedRight\arraybackslash}p{3.4cm}Y>{\RaggedRight\arraybackslash}p{4.6cm}}
\rowcolor{milcVino}\color{white}Dimensión & \color{white}\bfseries Mi respuesta (frame starter) & \color{white}\bfseries Algo que descubrí\\
1. Personal & Lo que más cambió en mí fue \linead{6.5cm} & \linead{4.2cm}\\
2. Emocional & La emoción más fuerte fue \linead{2.5cm} y la regulé \linead{3cm} & \linead{4.2cm}\\
3. Ciudadana & En democracia esto importa porque \linead{6cm} & \linead{4.2cm}\\
4. Local (Cartago) & En mi barrio sería útil para \linead{6.5cm} & \linead{4.2cm}\\
5. Intergeneracional & A mi abuela/abuelo le diría \linead{6.5cm} & \linead{4.2cm}\\
\end{tabularx}
\normalsize

\begin{infoband}{milcVino}
\textbf{Para la próxima sustentación.} Vuelve a esta tabla en seis meses. Verás que las respuestas cambian — no porque lo aprendido cambie, sino porque tú cambias. La evaluación liberadora es una conversación contigo a través del tiempo.
\end{infoband}


\puente{Y para terminar, tres voces sobre sustentar. Dussel sobre tomar la palabra pública, Marco Aurelio sobre la coherencia entre lo que dices y lo que hiciste, Floridi sobre la publicación como acto reputacional.}

\begin{titlebox}{milcGradeDark}
Triángulo de pensamiento · cierre filosófico
\end{titlebox}

\begin{softbox}{milcVino}{milcGris}{Dussel · lente del nosotros}
\textit{``Tomar la palabra pública es el primer acto de ciudadanía intelectual --- y se entrena en cada presentación, no se hereda.''}

Sustentar en público una pieza propia ante 30 personas es un acto político: te haces visible, asumes responsabilidad, soportas mirada. En Colombia, muchas personas pasan toda su vida escolar sin ese entrenamiento y luego no pueden defender sus ideas en el trabajo. Sustentar aquí, con apoyo del docente y los pares, es lo más cercano al gimnasio de la ciudadanía intelectual.

\textbf{¿Estoy acostumbrado a defender mis ideas en voz alta ante un grupo o tiendo a esconderme? ¿Qué significa esto para mi futuro adulto?}
\end{softbox}
\linead{\linewidth}\\[1mm]
\linead{\linewidth}

\begin{softbox}{milcMostaza}{milcGris}{Estoicismo · Marco Aurelio · cuidado interior}
\textit{``No es lo que dices al exponer; es la coherencia entre tu obra y tu palabra lo que se queda.''}

Puedes presentar una revista mediocre con palabras hermosas o una revista buena con palabras simples. La sabiduría estoica del sustentante: tus palabras solo amplifican lo que hizo tu obra. Si la obra está bien hecha, la sustentación la celebra. Si no, ninguna palabra la salva. Por eso \emph{el oficio precede a la palabra}.

\textbf{¿Mi revista está a la altura de lo que voy a decir de ella? Si no, ¿qué prefiero corregir: la obra o el discurso?}
\end{softbox}
\linead{\linewidth}\\[1mm]
\linead{\linewidth}

\begin{softbox}{milcTurquesa}{milcGris}{Floridi · lente de la infoesfera}
\textit{``Cada sustentación pública es un acto reputacional permanente en la infoesfera de tu vida.''}

Las personas que te escuchen hoy --- compañeros, profesores --- son la primera audiencia. Aprenden a leerte, a confiar o desconfiar de ti. Cada presentación construye reputación, lentamente, durante años. Sustentar bien hoy es inversión a largo plazo de tu identidad pública. No te la juegas por una nota --- te la juegas por quien eres ante los demás.

\textbf{¿Qué versión de mí va a ver mi audiencia hoy? ¿Esa versión es la que quiero proyectar en mi vida adulta?}
\end{softbox}
\linead{\linewidth}\\[1mm]
\linead{\linewidth}

\begin{titlebox}{milcVino}
Compromiso final
\end{titlebox}

\small
\begin{tabularx}{\textwidth}{>{\RaggedRight\arraybackslash}p{3.0cm}YYY}
\rowcolor{milcVino}\color{white}\bfseries Criterio & \color{white}\bfseries Lo logré & \color{white}\bfseries En proceso & \color{white}\bfseries Necesito apoyo\\
Comprensión del tema & Explico ideas centrales con mis palabras. & Reconozco ideas pero falta precisión. & Necesito repasar conceptos base.\\
Pensamiento computacional & Apliqué los 4 pilares al diseñar mi producto. & Apliqué algunos pilares. & Necesito releer la fase 2.\\
Aplicación y producto & Mi producto cumple los criterios y se puede revisar. & Producto funcional con ajustes pendientes. & Aún en construcción.\\
Reflexión 5 dimensiones & Respondí cada dimensión con honestidad. & Respondí algunas con detalle. & Mis respuestas son muy generales.\\
Triángulo de pensamiento & Conecté las 3 voces con mi trabajo real. & Conecté al menos una. & No relaciono las voces con mi proceso.\\
\end{tabularx}
\normalsize

\textbf{Hoy entendí que:}\\[1mm]
\linead{\linewidth}\\[1mm]
\linead{\linewidth}

\textbf{Mi compromiso para la próxima sesión es:}\\[1mm]
\linead{\linewidth}\\[1mm]
\linead{\linewidth}

\begin{softbox}{milcMostaza}{milcGris}{Cita de despedida · Marco Aurelio}
\textit{``El obstáculo en el camino se vuelve el camino. Lo que estorba al camino se vuelve camino.''}\\[1mm]
Cada error de hoy, bien observado, se convierte en pieza del camino que pisarás mañana.
\end{softbox}

\small
\begin{tabularx}{\textwidth}{>{\bfseries}p{3.6cm}Y}
\rowcolor{milcGradeDark!12}Nombre completo: & \linead{10cm}\\
Fecha: & \linead{4cm} \quad Curso/grupo: \linead{4cm}\\
Mi firma: & \linead{9cm}\\
\end{tabularx}
\normalsize

\begin{infoband}{milcGradeDark}
\textbf{Para el próximo periodo.} Pega o copia tu sustentación (estructura + feedback recibido) en la primera página de tu cuaderno del Periodo 3. Cada mes vuelve a leerla y anota: ¿estoy aplicando lo que me señalaron como debilidad? Esa práctica de \emph{volver al feedback} es lo que convierte una crítica en aprendizaje. \emph{Felicitaciones por cerrar el Periodo 2 con voz pública.}
\end{infoband}

\end{document}
